\chapterbackmatter{Supplementary Exercises}

\begin{enumerate}
  \SupplementaryExercise{1}{Cancellation Behind the Witten Index}
    {(Adapted from David Tong, \emph{Supersymmetric Field Theory},
    Chapter 3, Section 3.4)}
    {tong-sft-problems-2023-ch29}
  Let a Hermitian supercharge obey \(Q^2=H\) and anticommute with
  \((-1)^F\).  Prove that every positive-energy eigenspace has equal
  bosonic and fermionic dimensions, and determine which states can
  contribute to \(\operatorname{Tr}(-1)^F e^{-\beta H}\).

  \SupplementaryExercise{2}{A Vacuum Escaping to Infinity}
    {(Adapted from David Tong, \emph{Supersymmetric Field Theory},
    Chapter 3, Section 3.4)}
    {tong-sft-problems-2023-ch29}
  For one chiral superfield with
  \(f_\epsilon(X)=fX+\epsilon X^2/2\), find all supersymmetric vacua for
  \(\epsilon\ne0\), follow them as \(\epsilon\to0\), and explain precisely
  why index invariance does not forbid the limiting theory from breaking
  supersymmetry.

  \SupplementaryExercise{3}{Normalized Goldstino Direction}
    {(Adapted from David Tong, \emph{Supersymmetric Field Theory},
    Chapter 3, Section 3.4)}
    {tong-sft-problems-2023-ch29}
  In a general theory with chiral auxiliary expectation values
  \(\mathcal F_{n0}\) and gauge auxiliary expectation values \(D_{A0}\),
  construct the normalized massless fermion direction and express its
  normalization in terms of the vacuum energy density.

  \SupplementaryExercise{4}{Holomorphic Spurion Selection Rule}
    {(Adapted from David Tong, \emph{Supersymmetric Field Theory},
    Example Sheet 2, Question 1)}
    {tong-sft-problems-2023-ch29}
  Promote the coefficients in
  \(f(\Phi)=\mu_2\Phi^2+\mu_3\Phi^3\) to background chiral superfields.
  Assign spurionic phase and \(R\) charges, and use holomorphy plus the
  weak-coupling limit to rule out perturbative corrections to the
  Wilsonian superpotential.

  \SupplementaryExercise{5}{Supersymmetric One-Loop Gauge Coefficient}
    {(Adapted from David Tong, \emph{Supersymmetric Field Theory},
    Example Sheet 3, Question 1)}
    {tong-sft-problems-2023-ch29}
  Sum the gauge-boson, gaugino, Weyl-fermion, and complex-scalar
  contributions to derive \(b_0=3C_1-C_2\) for an \(N=1\)
  supersymmetric gauge theory.  Apply the result to \(SU(N_c)\) SQCD
  with \(N_f\) flavor pairs.

  \SupplementaryExercise{6}{Instanton Weight as a Chiral Parameter}
    {(Adapted from David Tong, \emph{Supersymmetric Field Theory},
    Example Sheet 3, Questions 4--5)}
    {tong-sft-problems-2023-ch29}
  Starting from the Euclidean instanton action and topological charge,
  show that a winding-\(\nu\) contribution carries
  \(\exp(2\pi i\nu\tau)\).  Explain why holomorphic Wilsonian
  \(\mathcal F\)-terms contain positive rather than negative powers of
  this quantity.

  \SupplementaryExercise{7}{Mass-Deformed ADS Vacua}
    {(Adapted from David Tong, \emph{Supersymmetric Field Theory},
    Example Sheet 4, Question 1)}
    {tong-sft-problems-2023-ch29}
  For \(SU(N_c)\) SQCD with \(N_f<N_c\), extremize
  \[
    f_{\rm eff}=(N_c-N_f)
    \left(\frac{\Lambda^{3N_c-N_f}}{\det M}\right)^{1/(N_c-N_f)}
    +\operatorname{Tr}(mM).
  \]
  Find \(M\), count the supersymmetric vacua, and derive the scale
  matching relation after all flavors are integrated out.

  \SupplementaryExercise{8}{Abelian Field-Strength \(\mathcal F\)-Term}
    {(Adapted from David Tong, \emph{Supersymmetric Field Theory},
    Example Sheet 2, Question 3)}
    {tong-sft-problems-2023-ch29}
  Expand the Abelian left-chiral field-strength superfield through
  quadratic order in \(\theta_L\).  Compute
  \([W_L^{\mathsf T}\epsilon W_L]_{\mathcal F}\), identifying separately
  the Maxwell, theta, gaugino, and auxiliary-field terms in Volume III
  conventions.

  \SupplementaryExercise{9}{Metric from the \(N=2\) Holomorphic Function}
    {(Adapted from Adel Bilal, \emph{Introduction to Supersymmetry},
    Sections 8.2 and 9.1)}
    {bilal-introduction-susy-2001-ch29}
  Given \(T(a)=h'(a)/(4\pi i)\) and
  \(K=\operatorname{Im}[a^*h(a)]/(4\pi)\), derive the scalar kinetic
  coefficient.  State the local positivity condition and explain why a
  logarithmic perturbative \(h(a)\) cannot remain valid near the origin.

  \SupplementaryExercise{10}{First-Order Abelian Dualization}
    {(Adapted from Adel Bilal, \emph{Introduction to Supersymmetry},
    Section 9.2.1)}
    {bilal-introduction-susy-2001-ch29}
  Treat \(W_L\) as unconstrained and impose its supersymmetric Bianchi
  identity with a real multiplier superfield.  Integrate out \(W_L\) and
  show that the dual scalar coordinate and coupling satisfy
  \(a_D=h(a)\) and \(T_D=1/(16\pi^2T)\).

  \SupplementaryExercise{11}{Generators of the Integral Duality Group}
    {(Adapted from Adel Bilal, \emph{Introduction to Supersymmetry},
    Section 9.2.2)}
    {bilal-introduction-susy-2001-ch29}
  Represent electric--magnetic interchange and an integral theta shift
  as matrices acting on \((a,a_D)^{\mathsf T}\).  Show that their
  products have integral entries and unit determinant, and determine the
  induced fractional-linear transformation of \(h'(a)\).

  \SupplementaryExercise{12}{Theta Shift and the Dyon Lattice}
    {(Adapted from Adel Bilal, \emph{Introduction to Supersymmetry},
    Sections 9.2.2--9.2.3)}
    {bilal-introduction-susy-2001-ch29}
  Use \(a_D\to a_D+ba\), with integral \(b\), to find the corresponding
  relabeling of electric and magnetic quantum numbers.  Verify that the
  BPS central charge and mass are unchanged and recover the Witten shift
  of a monopole's electric charge.

  \SupplementaryExercise{13}{Weak-Coupling Monodromy at Infinity}
    {(Adapted from Adel Bilal, \emph{Introduction to Supersymmetry},
    Sections 9.1.4 and 9.3.1)}
    {bilal-introduction-susy-2001-ch29}
  Starting from the one-loop asymptotic forms
  \(a\simeq\sqrt{2u}\) and
  \(a_D\simeq i\sqrt{2u}[\ln(2u/\Lambda^2)-2]/\pi\), carry \(u\)
  counterclockwise around infinity and derive its monodromy matrix.

  \SupplementaryExercise{14}{Monodromy of a Massless Monopole}
    {(Adapted from Adel Bilal, \emph{Introduction to Supersymmetry},
    Section 9.3.3)}
    {bilal-introduction-susy-2001-ch29}
  Near a point where a unit monopole becomes massless, take
  \(a_D\sim c(u-u_0)\) and
  \(a\sim a_0+(i/\pi)a_D\ln(a_D/\Lambda_0)\).
  Determine the monodromy acting on \((a,a_D)^{\mathsf T}\).

  \SupplementaryExercise{15}{Factorization of Seiberg--Witten Monodromies}
    {(Adapted from Adel Bilal, \emph{Introduction to Supersymmetry},
    Sections 9.3--9.4)}
    {bilal-introduction-susy-2001-ch29}
  Multiply the proposed monodromies at \(u=+1\) and \(u=-1\),
  \[
    M_+=\begin{pmatrix}1&-2\\0&1\end{pmatrix},
    \qquad
    M_-=\begin{pmatrix}3&-2\\2&-1\end{pmatrix}.
  \]
  Compare their ordered product with the weak-coupling monodromy and
  identify the massless charge vector at each singularity.

  \SupplementaryExercise{16}{Duality Covariance of the BPS Spectrum}
    {(Adapted from Adel Bilal, \emph{Introduction to Supersymmetry},
    Section 9.2.3)}
    {bilal-introduction-susy-2001-ch29}
  For \(Z=2\sqrt2\,i(na+ma_D)\), determine how the integer charge row
  \((n,m)\) must transform when
  \((a,a_D)^{\mathsf T}\to M(a,a_D)^{\mathsf T}\) with
  \(M\in SL(2,\mathbb Z)\).  Prove that the BPS mass is invariant.

  \SupplementaryExercise{17}{Wilsonian Versus 1PI Holomorphy}
    {(Adapted from K. Intriligator and N. Seiberg,
    \emph{Lectures on Supersymmetric Gauge Theories and
    Electric-Magnetic Duality}, Section 2.3)}
    {intriligator-seiberg-gauge-duality-1995-ch29}
  Explain why interacting massless particles can produce infrared
  holomorphic anomalies in a 1PI action but not in a Wilsonian action
  with a finite infrared cutoff.  Give a concrete criterion for when the
  two effective superpotentials may be identified.

  \SupplementaryExercise{18}{Anomaly-Free \(R\) Charge in SQCD}
    {(Adapted from K. Intriligator and N. Seiberg,
    \emph{Lectures on Supersymmetric Gauge Theories and
    Electric-Magnetic Duality}, Sections 3--4)}
    {intriligator-seiberg-gauge-duality-1995-ch29}
  In \(SU(N_c)\) SQCD with \(N_f\) flavor pairs, assign equal
  \(R\) charge to \(Q\) and \(\bar Q\).  Cancel the mixed
  \(SU(N_c)^2U(1)_R\) anomaly and find the \(R\) charges of the meson,
  baryons, and their fermionic components.

  \SupplementaryExercise{19}{Uniqueness of the ADS Superpotential}
    {(Adapted from K. Intriligator and N. Seiberg,
    \emph{Lectures on Supersymmetric Gauge Theories and
    Electric-Magnetic Duality}, Section 4.1)}
    {intriligator-seiberg-gauge-duality-1995-ch29}
  For \(N_f<N_c\), use flavor symmetry, the anomaly-free \(R\) symmetry,
  the anomalous axial symmetry, and mass dimension to derive the powers
  of \(\Lambda\) and \(\det M\) in the Affleck--Dine--Seiberg
  superpotential.

  \SupplementaryExercise{20}{Holomorphic Decoupling of One Flavor}
    {(Adapted from K. Intriligator and N. Seiberg,
    \emph{Lectures on Supersymmetric Gauge Theories and
    Electric-Magnetic Duality}, Section 4.1)}
    {intriligator-seiberg-gauge-duality-1995-ch29}
  Add \(mM_{N_fN_f}\) to the ADS superpotential, solve for the heavy
  meson, and show that the remaining light fields have the
  \(N_f-1\) flavor superpotential provided the two strong scales obey
  the correct matching relation.

  \SupplementaryExercise{21}{Gaugino Condensation and Vacuum Multiplicity}
    {(Adapted from K. Intriligator and N. Seiberg,
    \emph{Lectures on Supersymmetric Gauge Theories and
    Electric-Magnetic Duality}, Sections 1.2 and 4.1)}
    {intriligator-seiberg-gauge-duality-1995-ch29}
  Use the anomalous \(R\) symmetry of pure \(SU(N_c)\) super
  Yang--Mills theory to determine the unbroken discrete subgroup,
  enumerate the phases of the gaugino condensate, and relate the result
  to the Witten index.

  \SupplementaryExercise{22}{Mass Deformation of the Quantum Constraint}
    {(Adapted from K. Intriligator and N. Seiberg,
    \emph{Lectures on Supersymmetric Gauge Theories and
    Electric-Magnetic Duality}, Section 4.2)}
    {intriligator-seiberg-gauge-duality-1995-ch29}
  Impose \(\det M-B\bar B=\Lambda^{2N_c}\) for
  \(N_f=N_c\) with a chiral Lagrange multiplier.  Give one flavor a
  mass, integrate it and the multiplier out, and recover the
  \(N_f=N_c-1\) nonperturbative superpotential.

  \SupplementaryExercise{23}{Orthogonal-Group Holomorphic Dynamics}
    {(Adapted from K. Intriligator and N. Seiberg,
    \emph{Lectures on Supersymmetric Gauge Theories and
    Electric-Magnetic Duality}, Section 6)}
    {intriligator-seiberg-gauge-duality-1995-ch29}
  For \(SO(N_c)\) with \(N_f\) vector chiral superfields, derive the
  anomaly-free matter \(R\) charge and the one-loop coefficient.  For
  \(N_f<N_c-2\), determine the unique runaway superpotential built from
  the symmetric meson matrix.

  \SupplementaryExercise{24}{Broken Vacuum or Runaway?}
    {(Adapted from Erich Poppitz and Sandip P. Trivedi,
    \emph{Dynamical Supersymmetry Breaking}, Sections 1--2)}
    {poppitz-trivedi-dsb-1998-ch29}
  For a smooth positive Kahler metric, prove that a finite stationary
  point with no solution of \(f_i=0\) has positive energy and broken
  supersymmetry.  Explain why a direction with \(V\to0\) at infinite
  field values instead describes a runaway without a ground state.

  \SupplementaryExercise{25}{O'Raifeartaigh Pseudomodulus and Goldstino}
    {(Adapted from Erich Poppitz and Sandip P. Trivedi,
    \emph{Dynamical Supersymmetry Breaking}, Section 2)}
    {poppitz-trivedi-dsb-1998-ch29}
  Consider
  \(f=f_0X+\tfrac12hX\phi_1^2+m\phi_1\phi_2\), with
  \(m^2>\lvert hf_0\rvert\).  Find the classical vacuum manifold, its
  energy, the tree-level pseudomodulus, and the normalized goldstino
  direction.

  \SupplementaryExercise{26}{Radiative Lifting of a Pseudomodulus}
    {(Adapted from Erich Poppitz and Sandip P. Trivedi,
    \emph{Dynamical Supersymmetry Breaking}, Section 2)}
    {poppitz-trivedi-dsb-1998-ch29}
  For the O'Raifeartaigh model of the preceding exercise, write the
  Coleman--Weinberg supertrace formula.  Use the \(X\)-dependent mass
  matrices to determine the parametric size of the induced
  pseudomodulus squared mass and explain why the goldstino remains
  massless.

  \SupplementaryExercise{27}{A Generic \(R\)-Symmetry Counting Test}
    {(Adapted from Erich Poppitz and Sandip P. Trivedi,
    \emph{Dynamical Supersymmetry Breaking}, Section 2)}
    {poppitz-trivedi-dsb-1998-ch29}
  In a generic \(R\)-symmetric Wess--Zumino model, let \(N_X\) fields
  have \(R=2\) and \(N_Y\) fields have \(R=0\).  Analyze the
  \(\mathcal F\)-term equations at an \(R\)-symmetric point and state the
  condition under which supersymmetry is generically broken.

  \SupplementaryExercise{28}{The \(3\)-\(2\) Model Rank Obstruction}
    {(Adapted from Erich Poppitz and Sandip P. Trivedi,
    \emph{Dynamical Supersymmetry Breaking}, Section 4.1)}
    {poppitz-trivedi-dsb-1998-ch29}
  In the \(SU(3)\times SU(2)\) model with
  \(f_{\rm tree}=\lambda Q\bar D L\), combine the \(SU(3)\) ADS term with
  the \(L\)-field equation.  Show that no finite configuration is both
  \(\mathcal F\)-flat and nonsingular, and explain the role of
  \(D\)-terms in stabilizing the runaway.

  \SupplementaryExercise{29}{Quantum-Constraint Rank Condition}
    {(Adapted from Erich Poppitz and Sandip P. Trivedi,
    \emph{Dynamical Supersymmetry Breaking}, Section 4)}
    {poppitz-trivedi-dsb-1998-ch29}
  An \(Sp(N)\) theory with \(2N+2\) fundamentals has antisymmetric mesons
  satisfying \(\operatorname{Pf}M=\Lambda^{2N+2}\).  Couple singlets by
  \(f=\lambda S^{ij}M_{ij}\).  Prove that the singlet
  \(\mathcal F\)-terms are incompatible with the quantum constraint and
  identify the symmetry reason for the resulting goldstino.

  \SupplementaryExercise{30}{Current Normalization of the Goldstino Scale}
    {(Adapted from Erich Poppitz and Sandip P. Trivedi,
    \emph{Dynamical Supersymmetry Breaking}, Sections 1--2)}
    {poppitz-trivedi-dsb-1998-ch29}
  Parameterize the vacuum-to-goldstino matrix element of the conserved
  supercurrent by a constant \(F_G\).  Insert one-particle states into
  the supersymmetry algebra and derive both
  \(\rho_{\rm VAC}=\lvert F_G\rvert^2/2\) and the relation to chiral and
  gauge auxiliary expectation values.
\end{enumerate}
